\title{The Era of 1-bit LLMs: All Large Language Models are in 1.58 Bits}

\begin{abstract}
Emerging studies, notably BitNet~\cite{bitnet}, are ushering in a transformative phase for Large Language Models (LLMs) utilizing 1-bit precision. In this paper, we present a new 1-bit LLM, named \textbf{\text{BitNet b1.58}}, which constrains each model parameter (or weight) to the ternary set \{-1, 0, 1\}. This model achieves perplexity and downstream task results on par with standard full-precision (FP16 or BF16) Transformer-based LLMs of equivalent size and number of training tokens, all while offering major improvements in efficiency, including lower latency, reduced memory footprint, increased throughput, and diminished energy demands. Crucially, this 1.58-bit LLM proposes a novel scaling law and establishes new methodologies for the development of future LLMs that prioritize both performance and efficiency. Additionally, it lays the groundwork for a new computational framework and catalyzes innovation in hardware design specifically tailored to 1-bit LLMs.
\end{abstract}

\section{The Era of 1-bit LLMs}

Artificial intelligence research has recently witnessed substantial advances in the scale and effectiveness of Large Language Models (LLMs). Although these models excel in various natural language processing tasks, their rapid growth introduces significant obstacles regarding efficient deployment, as well as environmental and economic concerns stemming from high energy requirements.

A promising solution to these issues is post-training quantization, a method used to generate reduced-bitwidth models tailored for inference~\cite{smoothquant, gptq, quip, quip_sharp}. By lowering the numerical precision of both activations and weights, this approach drastically curtails memory use and computational complexity. The focus has progressively shifted from 16-bit formats towards models utilizing as few as 4 bits~\cite{gptq, awq}. Nevertheless, although post-training quantization has become a standard industry technique for LLMs, it remains sub-optimal in practice.

Emerging research into 1-bit model architectures, exemplified by BitNet~\cite{bitnet}, demonstrates substantial potential for reducing the computational expense of LLMs without sacrificing accuracy. Standard LLMs are typically implemented with 16-bit floating point representations (such as FP16 or BF16), and their dominant operations are matrix multiplications. This means that the computational overhead primarily arises from floating-point additions and multiplications. By contrast, BitNet’s matrix multiplications are based exclusively on integer addition, which dramatically reduces the associated energy demands. Given that power consumption is a core constraint for compute performance across most hardware accelerators, this efficiency also translates to increased computational speed.

Beyond pure computation, another bottleneck is the cost of moving model parameters from DRAM to on-chip accelerator memory (such as SRAM) during inference. Expanding SRAM capacity is one potential strategy to improve data throughput, but it results in significantly higher expenses compared to DRAM. In contrast, 1-bit LLMs exhibit a substantially reduced memory footprint in terms of both storage capacity and bandwidth requirements relative to their full-precision counterparts. This reduction in resource demand enables quicker and cheaper transfers of weights from DRAM, thereby facilitating faster and more economical inference.

In this paper, we present a noteworthy variant of the 1-bit LLM, termed \textbf{\text{BitNet b1.58}}, in which each model parameter is ternary and can assume values from the set \{-1, 0, 1\}. The original 1-bit BitNet architecture is expanded here by introducing the 0 value, leading to a representation that requires 1.58 bits in a binary scheme. \text{BitNet b1.58} maintains the foundational strengths of the original 1-bit BitNet: an innovative computation paradigm largely devoid of multiplications for matrix operations, resulting in exceptional efficiency. It preserves the same low energy profile and offers superior memory efficiency, throughput, and latency when benchmarked against FP16 LLMs. Importantly, \text{BitNet b1.58} introduces two further advantages. The inclusion of zero in its weight values enables explicit feature selection, thereby enhancing representational power and boosting 1-bit LLM performance. Empirically, our results indicate that, starting at model scales of 3B parameters, \text{BitNet b1.58} achieves parity with full-precision FP16 models across metrics such as perplexity and downstream task effectiveness, assuming consistent training configurations (model size, data volume, etc.).

\section{BitNet b1.58}

\text{BitNet b1.58} builds upon the BitNet framework, a Transformer variant where the standard \emph{nn.Linear} layers are substituted with \emph{BitLinear} layers. This model is trained from the beginning utilizing 1.58-bit precision for weights and 8-bit precision for activations. Relative to the original BitNet, several adjustments have been made, which are outlined as follows.

\paragraph{Quantization Function.} To enforce the weight values to be strictly -1, 0, or +1, an \emph{absmean} quantization approach is utilized. The method first normalizes the weight matrix by its mean absolute value, after which each entry is discretized to the closest integer among \{-1, 0, +1\}:

\begin{equation}
    \widetilde{W} = \mathrm{RoundClip}(\frac{W}{\gamma + \epsilon}, -1, 1),
\end{equation}

\begin{equation}
    \mathrm{RoundClip}(x, a, b) = \max(a, \min(b, \mathrm{round}(x))),
\end{equation}

\begin{equation}
    \gamma = \frac{1}{nm}\sum_{ij} |W_{ij}|.
\end{equation}

For activation quantization, the same strategy as BitNet is applied, with the exception that activations are not rescaled to the interval $[0, Q_b]$ before applying nonlinearities. Rather, for each token, activations are adjusted to fall within $[-Q_b, Q_b]$, eliminating zero-point quantization. This modification streamlines both implementation and system-level optimization, and, according to our empirical results, has only marginal impact on overall performance.

\paragraph{LLaMA-alike Components.}

The LLaMA architecture~\cite{llama, llama2} has established itself as a standard backbone for open-source large language models. In alignment with the open-source community, our design of \text{BitNet b1.58} integrates LLaMA-inspired components. Concretely, the model incorporates RMSNorm~\cite{rmsnorm}, SwiGLU~\cite{swiglu}, and rotary positional embedding~\cite{rope}, while omitting all bias terms. This facilitates seamless integration of \text{BitNet b1.58} into widely used open-source libraries such as Huggingface, vLLM~\cite{vllm}, and llama.cpp\footnote{\url{https://github.com/ggerganov/llama.cpp}}, requiring only minimal adaptation.

\section{Results}

We conducted a comparison between \text{BitNet b1.58} and our reproduced FP16 \text{LLaMA LLM} across different model scales. For consistency and fairness, both models were pre-trained on the RedPajama dataset~\cite{redpajama} using 100 billion tokens. Their zero-shot abilities were assessed on several language benchmarks, such as ARC-Easy~\cite{arc}, ARC-Challenge~\cite{arc}, Hellaswag~\cite{hellaswag}, Winogrande~\cite{winoGrande}, PIQA~\cite{piqa}, OpenbookQA~\cite{openbookqa}, and BoolQ~\cite{boolq}. In addition, validation perplexities were computed on WikiText2~\cite{wikitext2} and C4~\cite{c4} datasets.

To analyze efficiency, we measured runtime GPU memory usage and inference latency for both \text{LLaMA LLM} and \text{BitNet b1.58}. Measurements were performed using the FasterTransformer\footnote{\url{https://github.com/NVIDIA/FasterTransformer}} library, which provides high-performance LLM inference on GPU hardware. For \text{BitNet b1.58}, the 2-bit kernel from Ladder~\cite{ladder} was incorporated. The primary latency metric reported was time required per output token, reflecting inference cost.

Table~\ref{tab:ppl} presents the perplexity and associated costs for both \text{BitNet b1.58} and \text{LLaMA LLM}. The data indicate that \text{BitNet b1.58} achieves comparable perplexity to the FP16 \text{LLaMA LLM} at the 3B parameter scale, while being 2.71 times faster and requiring 3.55 times less GPU memory. Notably, the 3.9B parameter version of \text{BitNet b1.58} is 2.4 times faster, needs 3.32 times less GPU memory, and outperforms the 3B \text{LLaMA LLM}.

In Table~\ref{tab:zero-shot}, comprehensive zero-shot task results are reported. For evaluation, we used the standard approach from \emph{lm-evaluation-harness}\footnote{\url{https://github.com/EleutherAI/lm-evaluation-harness}}. The findings reveal that the gap in accuracy between \text{BitNet b1.58} and \text{LLaMA LLM} diminishes with larger model sizes. Importantly, \text{BitNet b1.58} achieves parity with the full-precision model starting at 3B parameters. Consistent with perplexity results, downstream task performance demonstrates that \text{BitNet b1.58} 3.9B not only outperforms \text{LLaMA LLM} 3B, but does so with markedly reduced memory footprint and latency. These results establish \text{BitNet b1.58} as a Pareto improvement relative to leading LLM architectures.

\paragraph{Memory and Latency}

To extend our analysis, we increased the model sizes to 7B, 13B, and 70B and assessed their associated computational costs. The patterns in latency and memory usage are displayed in Figure~\ref{fig:latency-memory-energy}, where we observe that acceleration benefits become more pronounced as the models grow larger. Notably, \text{BitNet b1.58} with 70B parameters achieves a 4.1$\times$ reduction in inference time compared to the \text{LLaMA LLM} reference model. This improvement stems from the increased computational burden of \emph{nn.Linear} with larger architectures. Memory usage follows a similar pattern, given that embeddings—kept in full precision—account for a diminishing fraction of the overall memory footprint as model size increases. All measurements of latency and memory were conducted using a 2-bit kernel, indicating potential for additional cost reductions through further kernel optimizations.

\paragraph{Energy}

We also analyze the energy required for arithmetic operations by both \text{BitNet b1.58} and \text{LLaMA LLM}, emphasizing the contribution of matrix multiplications since they dominate the energy profile in LLMs. As shown in Figure~\ref{fig:energy}, most of the energy expenditure in \text{BitNet b1.58} comes from INT8 additions, while \text{LLaMA LLM} relies heavily on both FP16 addition and multiplication operations. Referencing the energy model proposed in~\cite{energycost, pokebnn}, we find that \text{BitNet b1.58} offers a 71.4-fold decrease in the energy consumed by arithmetic operations during matrix multiplication on 7nm technology nodes. We additionally recorded total energy usage for processing inputs of 512 tokens. Our data indicates that as model capacity scales up, the relative energy efficiency advantage of \text{BitNet b1.58} over the FP16 \text{LLaMA LLM} baseline becomes even more significant. This trend can be attributed to the increasing share of \emph{nn.Linear} operations as models grow, alongside a reduced impact from non-linear components in larger architectures.

\paragraph{Throughput}

We evaluated the throughput of both \text{BitNet b1.58} and \text{LLaMA LLM} with 70B parameters on two A100 GPUs (80GB), employing pipeline parallelism~\cite{gpipe} to allow execution of the \text{LLaMA LLM} 70B model on this hardware. The batch size was incrementally increased up to the point at which GPU memory was fully utilized, maintaining a fixed sequence length of 512. As presented in Table~\ref{tab:throughput}, \text{BitNet b1.58} 70B accommodated a batch size up to 11 times larger than that of \text{LLaMA LLM}, translating to an 8.9-fold improvement in throughput.

\textbf{\text{BitNet b1.58} demonstrates a novel scaling relationship regarding the interplay between model efficacy and inference costs.} By referencing the data in Figures~\ref{fig:latency-memory-energy} and \ref{fig:energy}, one can draw the following size equivalences between 1.58-bit and conventional 16-bit models:

\begin{itemize}
    \item The 13B BitNet b1.58 outperforms a 3B FP16 LLM when considering latency, memory footprint, and energy usage.
    \item The 30B BitNet b1.58 delivers better latency, memory efficiency, and energy savings than the 7B FP16 LLM.
    \item The 70B BitNet b1.58 surpasses a 13B FP16 LLM on measures of latency, memory utilization, and energy expenditure.
\end{itemize}

\paragraph{Training with 2T Tokens}

Training token count plays a pivotal role in the development of LLMs. To assess the token scalability of \text{BitNet b1.58}, we trained this model variant with a dataset comprising 2T tokens, adhering to the same data curation protocol as StableLM-3B~\cite{StableLM-3B-4E1T}—a leading open-source 3B model. Evaluation was performed across a benchmark suite that includes Winogrande~\cite{winoGrande}, PIQA~\cite{piqa}, SciQ~\cite{sciq}, LAMBADA~\cite{lambada}, and ARC-easy~\cite{arc}, with zero-shot accuracy results presented in Table~\ref{tab:2t}. For benchmarks reporting both accuracy and normalized accuracy, their mean was computed. StableLM 3B results for the 2T-token training regime are sourced directly from its official documentation. Our experimental outcomes indicate that \text{BitNet b1.58} consistently surpasses StableLM 3B on all evaluated benchmarks, providing evidence that 1.58-bit LLMs possess robust generalization across tasks.

\section{Discussion and Future Work}

\textbf{1-bit Mixture-of-Experts (MoE) LLMs}

The Mixture-of-Experts (MoE) framework has established itself as an efficient solution to the computational demands of LLMs, substantially lowering the number of floating point operations (FLOPs) required. Nevertheless, deployment and scalability remain constrained by high memory usage and considerable inter-chip communication. Leveraging 1.58-bit LLMs provides viable remedies to these issues. The decreased memory requirements lead to fewer devices being necessary for MoE model deployment, and the lower memory overhead results in a substantial drop in the communication costs related to networked activation transfers. In optimal scenarios, where the entire MoE model can reside on a single chip, communication bottlenecks can be essentially eliminated.

\textbf{Native Support of Long Sequence in LLMs}

As LLMs advance, processing lengthy sequences has emerged as an essential capability. A key bottleneck for inference with extended context lengths is the memory consumed by KV caching mechanisms. \text{BitNet b1.58} moves toward native support for long-context processing by reducing activation precision from 16 bits to 8 bits, effectively doubling the feasible context window with the same hardware budget. Furthermore, these activations can be further compressed, in a lossless manner, to as little as 4 bits or even less, which could be particularly beneficial for 1.58-bit LLMs—an avenue we propose for future research.

\textbf{LLMs on Edge and Mobile}

Deploying 1.58-bit LLMs holds considerable promise for enhancing language model performance on resource-constrained edge and mobile platforms. These environments, which often operate under tight memory and computation budgets, typically struggle with running large-scale LLMs. The lowered memory and energy demands of 1.58-bit LLMs not only facilitate feasible deployment but also enable application scenarios previously unattainable on such devices, thereby expanding the functional range of edge and mobile systems. Additionally, since 1.58-bit LLMs operate efficiently on CPUs—the dominant processors in these platforms—\text{BitNet b1.58} can deliver improved efficiency and functionality directly on-device.

\textbf{New Hardware for 1-bit LLMs}

Recent advancements, exemplified by Groq\footnote{\url{https://groq.com/}}, indicate strong progress in purpose-built hardware architectures such as LPUs, specifically tailored to the requirements of LLM workloads. Building on these developments, we advocate for the creation of next-generation hardware and system designs explicitly optimized for 1-bit LLMs, thereby maximizing the advantages of the novel computational approach introduced with BitNet~\cite{bitnet}.

\end{document}